\subsection{Threats to Validity}
\label{sec:threatsvalidity}
% \todo{In this section, the most prominent threats to validity of your study should be discussed (the following categories are just an example suggestion). They have been moved with study design, as one should consider them before study execution.}

In the exploration of HPC in the Cloud one significant threat to validity is the potential presence of "implementation bias." This threat emerges when studies predominantly focus on successful implementations of HPC in the Cloud, possibly overlooking challenges, failures, or suboptimal instances.

For instance, researchers and industry professionals might be more inclined to publish and share success stories of HPC applications in the Cloud due to their positive outcomes. As a result, the literature may disproportionately highlight favorable scenarios, potentially giving a skewed perception of the actual challenges and limitations associated with implementing HPC solutions in cloud environments.

To address this threat, it is crucial for the literature review to actively seek out and include studies that discuss unsuccessful or challenging implementations of HPC in the Cloud. Incorporating a diverse range of perspectives and outcomes ensures a more comprehensive and realistic understanding of the successes and limitations within this dynamic intersection of High-Performance Computing and Cloud technology.

\subsubsection{External Validity}
In the context of external validity for HPC in the Cloud, one potential concern is the generalizability of findings across different cloud service providers and HPC architectures. Variability in infrastructure, performance capabilities, and resource allocation methods among different cloud platforms may impact the external validity of research findings. To enhance external validity, it is essential to consider a diverse range of cloud environments and HPC configurations in the literature review.

\subsubsection{Internal Validity}
Internal validity is crucial for ensuring the accuracy and reliability of findings within the realm of HPC in the Cloud. Potential threats may arise from variations in experiment design, such as differences in workload characteristics, benchmarking methodologies, and measurement tools. Researchers should carefully evaluate and account for these factors to maintain internal validity, providing a solid foundation for drawing accurate conclusions about the performance and efficiency of HPC applications in cloud environments.

\subsubsection{Construct Validity}
Construct validity in the study of HPC in the Cloud involves ensuring that the chosen measures and metrics accurately capture the intended aspects of performance, scalability, and efficiency. For example, many studies refer to large-scale applications while testing on small-scale infrastructures to limit costs, which is not contributing effectively to this literature review. Challenges may arise from the use of diverse performance metrics across studies, making direct comparisons difficult. Therefore, ensuring a standardized set of metrics and benchmarks relevant to HPC workloads in the Cloud can enhance construct validity, allowing for meaningful comparisons and a more cohesive understanding of the field.

\subsubsection{Conclusion Validity}
In the context of HPC in the Cloud pertains to the accuracy and appropriateness of inferences drawn from the reviewed literature. Potential threats include incomplete reporting of results, selective publication, or biased interpretation of findings. This is an important factor to consider given that the contributions listed in \autoref{sec:research_goal} might be compromised. To enhance conclusion validity, it is crucial to critically evaluate the methodologies and interpretations presented in the literature, ensuring that the conclusions drawn accurately reflect the nuances and complexities of implementing HPC in diverse cloud computing environments.
